% Options for packages loaded elsewhere
% Options for packages loaded elsewhere
\PassOptionsToPackage{unicode}{hyperref}
\PassOptionsToPackage{hyphens}{url}
\PassOptionsToPackage{dvipsnames,svgnames,x11names}{xcolor}
%
\documentclass[
  letterpaper,
  DIV=11,
  numbers=noendperiod]{scrartcl}
\usepackage{xcolor}
\usepackage[top=30mm,left=20mm,heightrounded]{geometry}
\usepackage{amsmath,amssymb}
\setcounter{secnumdepth}{-\maxdimen} % remove section numbering
\usepackage{iftex}
\ifPDFTeX
  \usepackage[T1]{fontenc}
  \usepackage[utf8]{inputenc}
  \usepackage{textcomp} % provide euro and other symbols
\else % if luatex or xetex
  \usepackage{unicode-math} % this also loads fontspec
  \defaultfontfeatures{Scale=MatchLowercase}
  \defaultfontfeatures[\rmfamily]{Ligatures=TeX,Scale=1}
\fi
\usepackage{lmodern}
\ifPDFTeX\else
  % xetex/luatex font selection
\fi
% Use upquote if available, for straight quotes in verbatim environments
\IfFileExists{upquote.sty}{\usepackage{upquote}}{}
\IfFileExists{microtype.sty}{% use microtype if available
  \usepackage[]{microtype}
  \UseMicrotypeSet[protrusion]{basicmath} % disable protrusion for tt fonts
}{}
\makeatletter
\@ifundefined{KOMAClassName}{% if non-KOMA class
  \IfFileExists{parskip.sty}{%
    \usepackage{parskip}
  }{% else
    \setlength{\parindent}{0pt}
    \setlength{\parskip}{6pt plus 2pt minus 1pt}}
}{% if KOMA class
  \KOMAoptions{parskip=half}}
\makeatother
% Make \paragraph and \subparagraph free-standing
\makeatletter
\ifx\paragraph\undefined\else
  \let\oldparagraph\paragraph
  \renewcommand{\paragraph}{
    \@ifstar
      \xxxParagraphStar
      \xxxParagraphNoStar
  }
  \newcommand{\xxxParagraphStar}[1]{\oldparagraph*{#1}\mbox{}}
  \newcommand{\xxxParagraphNoStar}[1]{\oldparagraph{#1}\mbox{}}
\fi
\ifx\subparagraph\undefined\else
  \let\oldsubparagraph\subparagraph
  \renewcommand{\subparagraph}{
    \@ifstar
      \xxxSubParagraphStar
      \xxxSubParagraphNoStar
  }
  \newcommand{\xxxSubParagraphStar}[1]{\oldsubparagraph*{#1}\mbox{}}
  \newcommand{\xxxSubParagraphNoStar}[1]{\oldsubparagraph{#1}\mbox{}}
\fi
\makeatother

\usepackage{color}
\usepackage{fancyvrb}
\newcommand{\VerbBar}{|}
\newcommand{\VERB}{\Verb[commandchars=\\\{\}]}
\DefineVerbatimEnvironment{Highlighting}{Verbatim}{commandchars=\\\{\}}
% Add ',fontsize=\small' for more characters per line
\usepackage{framed}
\definecolor{shadecolor}{RGB}{241,243,245}
\newenvironment{Shaded}{\begin{snugshade}}{\end{snugshade}}
\newcommand{\AlertTok}[1]{\textcolor[rgb]{0.68,0.00,0.00}{#1}}
\newcommand{\AnnotationTok}[1]{\textcolor[rgb]{0.37,0.37,0.37}{#1}}
\newcommand{\AttributeTok}[1]{\textcolor[rgb]{0.40,0.45,0.13}{#1}}
\newcommand{\BaseNTok}[1]{\textcolor[rgb]{0.68,0.00,0.00}{#1}}
\newcommand{\BuiltInTok}[1]{\textcolor[rgb]{0.00,0.23,0.31}{#1}}
\newcommand{\CharTok}[1]{\textcolor[rgb]{0.13,0.47,0.30}{#1}}
\newcommand{\CommentTok}[1]{\textcolor[rgb]{0.37,0.37,0.37}{#1}}
\newcommand{\CommentVarTok}[1]{\textcolor[rgb]{0.37,0.37,0.37}{\textit{#1}}}
\newcommand{\ConstantTok}[1]{\textcolor[rgb]{0.56,0.35,0.01}{#1}}
\newcommand{\ControlFlowTok}[1]{\textcolor[rgb]{0.00,0.23,0.31}{\textbf{#1}}}
\newcommand{\DataTypeTok}[1]{\textcolor[rgb]{0.68,0.00,0.00}{#1}}
\newcommand{\DecValTok}[1]{\textcolor[rgb]{0.68,0.00,0.00}{#1}}
\newcommand{\DocumentationTok}[1]{\textcolor[rgb]{0.37,0.37,0.37}{\textit{#1}}}
\newcommand{\ErrorTok}[1]{\textcolor[rgb]{0.68,0.00,0.00}{#1}}
\newcommand{\ExtensionTok}[1]{\textcolor[rgb]{0.00,0.23,0.31}{#1}}
\newcommand{\FloatTok}[1]{\textcolor[rgb]{0.68,0.00,0.00}{#1}}
\newcommand{\FunctionTok}[1]{\textcolor[rgb]{0.28,0.35,0.67}{#1}}
\newcommand{\ImportTok}[1]{\textcolor[rgb]{0.00,0.46,0.62}{#1}}
\newcommand{\InformationTok}[1]{\textcolor[rgb]{0.37,0.37,0.37}{#1}}
\newcommand{\KeywordTok}[1]{\textcolor[rgb]{0.00,0.23,0.31}{\textbf{#1}}}
\newcommand{\NormalTok}[1]{\textcolor[rgb]{0.00,0.23,0.31}{#1}}
\newcommand{\OperatorTok}[1]{\textcolor[rgb]{0.37,0.37,0.37}{#1}}
\newcommand{\OtherTok}[1]{\textcolor[rgb]{0.00,0.23,0.31}{#1}}
\newcommand{\PreprocessorTok}[1]{\textcolor[rgb]{0.68,0.00,0.00}{#1}}
\newcommand{\RegionMarkerTok}[1]{\textcolor[rgb]{0.00,0.23,0.31}{#1}}
\newcommand{\SpecialCharTok}[1]{\textcolor[rgb]{0.37,0.37,0.37}{#1}}
\newcommand{\SpecialStringTok}[1]{\textcolor[rgb]{0.13,0.47,0.30}{#1}}
\newcommand{\StringTok}[1]{\textcolor[rgb]{0.13,0.47,0.30}{#1}}
\newcommand{\VariableTok}[1]{\textcolor[rgb]{0.07,0.07,0.07}{#1}}
\newcommand{\VerbatimStringTok}[1]{\textcolor[rgb]{0.13,0.47,0.30}{#1}}
\newcommand{\WarningTok}[1]{\textcolor[rgb]{0.37,0.37,0.37}{\textit{#1}}}

\usepackage{longtable,booktabs,array}
\newcounter{none} % for unnumbered tables
\usepackage{calc} % for calculating minipage widths
% Correct order of tables after \paragraph or \subparagraph
\usepackage{etoolbox}
\makeatletter
\patchcmd\longtable{\par}{\if@noskipsec\mbox{}\fi\par}{}{}
\makeatother
% Allow footnotes in longtable head/foot
\IfFileExists{footnotehyper.sty}{\usepackage{footnotehyper}}{\usepackage{footnote}}
\makesavenoteenv{longtable}
\usepackage{graphicx}
\makeatletter
\newsavebox\pandoc@box
\newcommand*\pandocbounded[1]{% scales image to fit in text height/width
  \sbox\pandoc@box{#1}%
  \Gscale@div\@tempa{\textheight}{\dimexpr\ht\pandoc@box+\dp\pandoc@box\relax}%
  \Gscale@div\@tempb{\linewidth}{\wd\pandoc@box}%
  \ifdim\@tempb\p@<\@tempa\p@\let\@tempa\@tempb\fi% select the smaller of both
  \ifdim\@tempa\p@<\p@\scalebox{\@tempa}{\usebox\pandoc@box}%
  \else\usebox{\pandoc@box}%
  \fi%
}
% Set default figure placement to htbp
\def\fps@figure{htbp}
\makeatother





\setlength{\emergencystretch}{3em} % prevent overfull lines

\providecommand{\tightlist}{%
  \setlength{\itemsep}{0pt}\setlength{\parskip}{0pt}}



 


\usepackage{booktabs}
\usepackage{longtable}
\usepackage{array}
\usepackage{multirow}
\usepackage{wrapfig}
\usepackage{float}
\usepackage{colortbl}
\usepackage{pdflscape}
\usepackage{tabu}
\usepackage{threeparttable}
\usepackage{threeparttablex}
\usepackage[normalem]{ulem}
\usepackage{makecell}
\usepackage{xcolor}
\KOMAoption{captions}{tableheading}
\makeatletter
\@ifpackageloaded{caption}{}{\usepackage{caption}}
\AtBeginDocument{%
\ifdefined\contentsname
  \renewcommand*\contentsname{Table of contents}
\else
  \newcommand\contentsname{Table of contents}
\fi
\ifdefined\listfigurename
  \renewcommand*\listfigurename{List of Figures}
\else
  \newcommand\listfigurename{List of Figures}
\fi
\ifdefined\listtablename
  \renewcommand*\listtablename{List of Tables}
\else
  \newcommand\listtablename{List of Tables}
\fi
\ifdefined\figurename
  \renewcommand*\figurename{Figure}
\else
  \newcommand\figurename{Figure}
\fi
\ifdefined\tablename
  \renewcommand*\tablename{Table}
\else
  \newcommand\tablename{Table}
\fi
}
\@ifpackageloaded{float}{}{\usepackage{float}}
\floatstyle{ruled}
\@ifundefined{c@chapter}{\newfloat{codelisting}{h}{lop}}{\newfloat{codelisting}{h}{lop}[chapter]}
\floatname{codelisting}{Listing}
\newcommand*\listoflistings{\listof{codelisting}{List of Listings}}
\makeatother
\makeatletter
\makeatother
\makeatletter
\@ifpackageloaded{caption}{}{\usepackage{caption}}
\@ifpackageloaded{subcaption}{}{\usepackage{subcaption}}
\makeatother
\usepackage{bookmark}
\IfFileExists{xurl.sty}{\usepackage{xurl}}{} % add URL line breaks if available
\urlstyle{same}
% fallback for those not using the hyperref driver hyperxmp:
\makeatletter
\@ifundefined{xmpquote}{\newcommand{\xmpquote}[1]{#1}}{}
\makeatother
\hypersetup{
  pdftitle={Computation for Health Data - Week 3},
  pdfauthor={Vicki Yiannakis},
  colorlinks=true,
  linkcolor={blue},
  filecolor={Maroon},
  citecolor={Blue},
  urlcolor={Blue},
  pdfcreator={LaTeX via pandoc}}


\title{Computation for Health Data - Week 3}
\author{Vicki Yiannakis}
\date{}
\begin{document}
\maketitle


First we will read in the maternal\_mortality CSV

\begin{Shaded}
\begin{Highlighting}[]
\NormalTok{matmor0 }\OtherTok{\textless{}{-}} \FunctionTok{read.csv}\NormalTok{(}\StringTok{"../data/raw/maternal\_mortality.csv"}\NormalTok{, }\AttributeTok{header =} \ConstantTok{TRUE}\NormalTok{)}
\FunctionTok{glimpse}\NormalTok{(matmor0)}
\end{Highlighting}
\end{Shaded}

\begin{verbatim}
Rows: 186
Columns: 22
$ iso       <chr> "AFG", "AGO", "ALB", "AND", "ARE", "ARG", "ARM", "ATG", "AUS~
$ indicator <chr> "Maternal mortality ratio (modeled estimate, per 100,000 liv~
$ X2000     <int> 1450, 827, 23, NA, 6, 66, 43, 44, 7, 6, 47, 1010, 8, 520, 51~
$ X2001     <int> 1390, 766, 23, NA, 6, 67, 42, 44, 6, 6, 44, 956, 8, 516, 501~
$ X2002     <int> 1300, 690, 21, NA, 5, 65, 39, 43, 6, 6, 43, 925, 7, 511, 486~
$ X2003     <int> 1240, 628, 21, NA, 5, 65, 38, 44, 6, 6, 44, 890, 7, 510, 471~
$ X2004     <int> 1180, 574, 18, NA, 5, 61, 36, 41, 5, 5, 39, 844, 7, 505, 454~
$ X2005     <int> 1140, 519, 22, NA, 5, 59, 35, 40, 5, 6, 42, 814, 7, 500, 437~
$ X2006     <int> 1120, 473, 18, NA, 4, 57, 36, 46, 5, 5, 35, 785, 7, 493, 422~
$ X2007     <int> 1090, 431, 19, NA, 4, 56, 32, 48, 5, 5, 34, 756, 7, 486, 410~
$ X2008     <int> 1030, 395, 20, NA, 4, 53, 36, 50, 5, 5, 32, 733, 7, 480, 401~
$ X2009     <int> 993, 359, 20, NA, 4, 56, 32, 45, 5, 5, 32, 698, 6, 471, 393,~
$ X2010     <int> 954, 326, 21, NA, 4, 51, 32, 44, 5, 5, 31, 665, 6, 464, 385,~
$ X2011     <int> 905, 300, 22, NA, 4, 48, 30, 43, 6, 5, 30, 635, 6, 458, 377,~
$ X2012     <int> 858, 281, 17, NA, 4, 47, 30, 44, 6, 5, 29, 608, 6, 450, 369,~
$ X2013     <int> 810, 269, 16, NA, 3, 44, 26, 43, 6, 5, 28, 591, 6, 441, 362,~
$ X2014     <int> 786, 258, 16, NA, 3, 42, 27, 42, 6, 5, 28, 576, 6, 432, 353,~
$ X2015     <int> 701, 251, 15, NA, 3, 41, 28, 43, 6, 5, 27, 568, 5, 421, 343,~
$ X2016     <int> 673, 246, 16, NA, 3, 40, 26, 43, 6, 5, 26, 558, 5, 408, 331,~
$ X2017     <int> 638, 241, 15, NA, 3, 39, 26, 42, 6, 5, 26, 548, 5, 397, 320,~
$ X2018     <lgl> NA, NA, NA, NA, NA, NA, NA, NA, NA, NA, NA, NA, NA, NA, NA, ~
$ X2019     <lgl> NA, NA, NA, NA, NA, NA, NA, NA, NA, NA, NA, NA, NA, NA, NA, ~
\end{verbatim}

We want to create a function that uses the below code pivots this to
long form

\begin{Shaded}
\begin{Highlighting}[]
\CommentTok{\# Change wide to long format}
\NormalTok{matmor\_long }\OtherTok{\textless{}{-}}\NormalTok{ matmor0 }\SpecialCharTok{|\textgreater{}}
  \FunctionTok{pivot\_longer}\NormalTok{(}\AttributeTok{cols =} \FunctionTok{starts\_with}\NormalTok{(}\StringTok{"X"}\NormalTok{), }\CommentTok{\#selecting all cols that start with x }
                 \AttributeTok{names\_to =} \StringTok{"year"}\NormalTok{, }\CommentTok{\#creating new variable called Year }
                 \AttributeTok{names\_prefix =} \StringTok{"X"}\NormalTok{, }\CommentTok{\# removes X from year column (so only numerical shows)}
                 \AttributeTok{values\_to =} \StringTok{"maternal\_mortality"}\NormalTok{) }\SpecialCharTok{|\textgreater{}} \CommentTok{\#values of maternal mortality }
  \FunctionTok{mutate}\NormalTok{(}\AttributeTok{year =} \FunctionTok{as.numeric}\NormalTok{(year)) }\SpecialCharTok{|\textgreater{}} \CommentTok{\# change year to numeric As it was stored as a character var}
  \FunctionTok{select}\NormalTok{(iso, year, maternal\_mortality) }\CommentTok{\#output data }
\end{Highlighting}
\end{Shaded}

NOTES: names\_prefix = A regular expression used to remove matching text
from the start of each variable name.

To create a function, we need to figure out our arguments

\begin{Shaded}
\begin{Highlighting}[]
\NormalTok{long\_form }\OtherTok{\textless{}{-}} \ControlFlowTok{function}\NormalTok{(x, varname) \{}
\NormalTok{    x}\SpecialCharTok{|\textgreater{}}
    \FunctionTok{pivot\_longer}\NormalTok{(}\AttributeTok{cols =} \FunctionTok{starts\_with}\NormalTok{(}\StringTok{"X"}\NormalTok{), }\CommentTok{\#selecting all cols that start with x }
                 \AttributeTok{names\_to =} \StringTok{"year"}\NormalTok{, }\CommentTok{\#creating new variable called Year }
                 \AttributeTok{names\_prefix =} \StringTok{"X"}\NormalTok{, }\CommentTok{\# removes X from year column (so only numerical shows)}
                 \AttributeTok{values\_to =}\NormalTok{ varname)}\SpecialCharTok{|\textgreater{}}
        \FunctionTok{mutate}\NormalTok{(}\AttributeTok{year =} \FunctionTok{as.numeric}\NormalTok{(year)) }\SpecialCharTok{|\textgreater{}}
        \FunctionTok{select}\NormalTok{(iso,year, }\FunctionTok{all\_of}\NormalTok{(varname)) }\CommentTok{\#values of maternal mortality }
\NormalTok{\} }

\CommentTok{\#now we can try this in the original mortality data}
\NormalTok{long\_maternal }\OtherTok{\textless{}{-}} \FunctionTok{long\_form}\NormalTok{(matmor0, }\StringTok{"maternal\_mortaltiy"}\NormalTok{)}
\end{Highlighting}
\end{Shaded}

Now we will try with the other datasets!

\begin{Shaded}
\begin{Highlighting}[]
\NormalTok{neonatal }\OtherTok{\textless{}{-}} \FunctionTok{read.csv}\NormalTok{(}\StringTok{"../data/raw/neonatal\_mortality.csv"}\NormalTok{, }\AttributeTok{header =} \ConstantTok{TRUE}\NormalTok{)}
\NormalTok{neonatal\_long }\OtherTok{\textless{}{-}} \FunctionTok{long\_form}\NormalTok{(neonatal, }\StringTok{"Neonatal\_Mortality"}\NormalTok{)}


\NormalTok{infant }\OtherTok{\textless{}{-}} \FunctionTok{read.csv}\NormalTok{(}\StringTok{"../data/raw/infant\_mortality.csv"}\NormalTok{, }\AttributeTok{header=}\ConstantTok{TRUE}\NormalTok{)}
\NormalTok{infant\_long }\OtherTok{\textless{}{-}} \FunctionTok{long\_form}\NormalTok{(infant, }\StringTok{"Infant\_Mortality"}\NormalTok{)}

\NormalTok{under5 }\OtherTok{\textless{}{-}} \FunctionTok{read.csv}\NormalTok{(}\StringTok{"../data/raw/under5\_mortality.csv"}\NormalTok{, }\AttributeTok{header=}\ConstantTok{TRUE}\NormalTok{)}
\NormalTok{under5\_long }\OtherTok{\textless{}{-}} \FunctionTok{long\_form}\NormalTok{(under5, }\StringTok{"Under5\_Mortality"}\NormalTok{)}
\end{Highlighting}
\end{Shaded}

\begin{enumerate}
\def\labelenumi{\arabic{enumi}.}
\setcounter{enumi}{1}
\tightlist
\item
  Prepare disaster data
\end{enumerate}

\begin{Shaded}
\begin{Highlighting}[]
\NormalTok{disaster }\OtherTok{\textless{}{-}} \FunctionTok{read.csv}\NormalTok{(}\StringTok{"../data/raw/disaster.csv"}\NormalTok{, }\AttributeTok{header=}\ConstantTok{TRUE}\NormalTok{)}
\NormalTok{disaster1 }\OtherTok{\textless{}{-}} \FunctionTok{clean\_names}\NormalTok{(disaster)}
\end{Highlighting}
\end{Shaded}

Now we will filter to only include the variables we need Make sure to
check variable types!

\begin{Shaded}
\begin{Highlighting}[]
\FunctionTok{class}\NormalTok{(disaster1}\SpecialCharTok{$}\NormalTok{year)}
\end{Highlighting}
\end{Shaded}

\begin{verbatim}
[1] "integer"
\end{verbatim}

\begin{Shaded}
\begin{Highlighting}[]
\FunctionTok{class}\NormalTok{(disaster1}\SpecialCharTok{$}\NormalTok{disaster\_type)}
\end{Highlighting}
\end{Shaded}

\begin{verbatim}
[1] "character"
\end{verbatim}

\begin{Shaded}
\begin{Highlighting}[]
\NormalTok{disaster\_filter }\OtherTok{\textless{}{-}}\NormalTok{ disaster1 }\SpecialCharTok{\%\textgreater{}\%} \FunctionTok{filter}\NormalTok{(year }\SpecialCharTok{\textgreater{}=} \DecValTok{2009} \SpecialCharTok{\&}\NormalTok{ year }\SpecialCharTok{\textless{}=} \DecValTok{2019}\NormalTok{,}
\NormalTok{                                        disaster\_type }\SpecialCharTok{==} \StringTok{"Earthquake"} \SpecialCharTok{|} 
\NormalTok{                                            disaster\_type }\SpecialCharTok{==} \StringTok{"Drought"}\NormalTok{) }\SpecialCharTok{\%\textgreater{}\%} 
                    \FunctionTok{select}\NormalTok{(year, iso, disaster\_type)}

\FunctionTok{kable}\NormalTok{(disaster\_filter}\SpecialCharTok{|\textgreater{}} \FunctionTok{slice}\NormalTok{(}\DecValTok{1}\SpecialCharTok{:}\DecValTok{10}\NormalTok{))                                }
\end{Highlighting}
\end{Shaded}

{\def\LTcaptype{none} % do not increment counter
\begin{longtable}[]{@{}rll@{}}
\toprule\noalign{}
year & iso & disaster\_type \\
\midrule\noalign{}
\endhead
\bottomrule\noalign{}
\endlastfoot
2009 & CRI & Earthquake \\
2009 & HND & Earthquake \\
2009 & AFG & Earthquake \\
2009 & CHN & Earthquake \\
2009 & IDN & Earthquake \\
2009 & IDN & Earthquake \\
2009 & ITA & Earthquake \\
2009 & JPN & Earthquake \\
2009 & IDN & Earthquake \\
2009 & IDN & Earthquake \\
\end{longtable}
}

Create a dummy variable!

\begin{Shaded}
\begin{Highlighting}[]
\NormalTok{disaster\_filter2 }\OtherTok{\textless{}{-}}\NormalTok{ disaster\_filter }\SpecialCharTok{\%\textgreater{}\%} 
    \FunctionTok{mutate}\NormalTok{(}\AttributeTok{earthquake =} \FunctionTok{if\_else}\NormalTok{(disaster\_type}\SpecialCharTok{==}\StringTok{"Earthquake"}\NormalTok{, }\DecValTok{1}\NormalTok{, }\DecValTok{0}\NormalTok{),}
           \AttributeTok{drought =} \FunctionTok{if\_else}\NormalTok{(disaster\_type}\SpecialCharTok{==}\StringTok{"Drought"}\NormalTok{, }\DecValTok{1}\NormalTok{, }\DecValTok{0}\NormalTok{)) }\SpecialCharTok{\%\textgreater{}\%}
                    \FunctionTok{mutate}\NormalTok{((}\FunctionTok{across}\NormalTok{(}\FunctionTok{where}\NormalTok{(is.numeric), as.integer)))}
                                                

\FunctionTok{kable}\NormalTok{(disaster\_filter2}\SpecialCharTok{|\textgreater{}} \FunctionTok{slice}\NormalTok{(}\DecValTok{1}\SpecialCharTok{:}\DecValTok{10}\NormalTok{))}
\end{Highlighting}
\end{Shaded}

{\def\LTcaptype{none} % do not increment counter
\begin{longtable}[]{@{}rllrr@{}}
\toprule\noalign{}
year & iso & disaster\_type & earthquake & drought \\
\midrule\noalign{}
\endhead
\bottomrule\noalign{}
\endlastfoot
2009 & CRI & Earthquake & 1 & 0 \\
2009 & HND & Earthquake & 1 & 0 \\
2009 & AFG & Earthquake & 1 & 0 \\
2009 & CHN & Earthquake & 1 & 0 \\
2009 & IDN & Earthquake & 1 & 0 \\
2009 & IDN & Earthquake & 1 & 0 \\
2009 & ITA & Earthquake & 1 & 0 \\
2009 & JPN & Earthquake & 1 & 0 \\
2009 & IDN & Earthquake & 1 & 0 \\
2009 & IDN & Earthquake & 1 & 0 \\
\end{longtable}
}

\begin{enumerate}
\def\labelenumi{\arabic{enumi}.}
\setcounter{enumi}{2}
\tightlist
\item
  Prepare the disaster data file
\end{enumerate}

We are creating a binary variable to ascertain whether they had
armed\_conflict in that year

\begin{Shaded}
\begin{Highlighting}[]
\NormalTok{conflict }\OtherTok{\textless{}{-}} \FunctionTok{read.csv}\NormalTok{(}\StringTok{"C:/Users/iamvi/OneDrive {-} University of Toronto/Year 2/Computation for Health Data/armed\_conflict/data/raw/conflict.csv"}\NormalTok{, }\AttributeTok{header=}\ConstantTok{TRUE}\NormalTok{)}

\CommentTok{\#we will need to create a binary\_conflict variable for each year 1999 to 2018 (19 total)}
\CommentTok{\#first we need to filter and sum up the number of conflicts each year}

\NormalTok{conflict1 }\OtherTok{\textless{}{-}}\NormalTok{ conflict }\SpecialCharTok{\%\textgreater{}\%} \FunctionTok{group\_by}\NormalTok{(iso,year)}\SpecialCharTok{\%\textgreater{}\%} \FunctionTok{summarise}\NormalTok{(}\AttributeTok{total\_conflicts=}\FunctionTok{sum}\NormalTok{(best))}
\end{Highlighting}
\end{Shaded}

\begin{verbatim}
`summarise()` has regrouped the output.
i Summaries were computed grouped by iso and year.
i Output is grouped by iso.
i Use `summarise(.groups = "drop_last")` to silence this message.
i Use `summarise(.by = c(iso, year))` for per-operation grouping
  (`?dplyr::dplyr_by`) instead.
\end{verbatim}

\begin{Shaded}
\begin{Highlighting}[]
\CommentTok{\#now we are going to create a binary variable in which year is lagged }


\NormalTok{conflict1 }\OtherTok{\textless{}{-}}\NormalTok{ conflict }\SpecialCharTok{\%\textgreater{}\%} \FunctionTok{group\_by}\NormalTok{(iso,year)}\SpecialCharTok{\%\textgreater{}\%} \FunctionTok{summarise}\NormalTok{(}\AttributeTok{total\_conflicts=}\FunctionTok{sum}\NormalTok{(best, }\AttributeTok{na.rm=}\ConstantTok{TRUE}\NormalTok{))}\SpecialCharTok{\%\textgreater{}\%} \FunctionTok{mutate}\NormalTok{(}\AttributeTok{binary\_conflict =} \FunctionTok{if\_else}\NormalTok{(}\FunctionTok{lead}\NormalTok{(total\_conflicts)}\SpecialCharTok{\textgreater{}}\DecValTok{0}\NormalTok{, }\DecValTok{1}\NormalTok{, }\DecValTok{0}\NormalTok{))}
\end{Highlighting}
\end{Shaded}

\begin{verbatim}
`summarise()` has regrouped the output.
i Summaries were computed grouped by iso and year.
i Output is grouped by iso.
i Use `summarise(.groups = "drop_last")` to silence this message.
i Use `summarise(.by = c(iso, year))` for per-operation grouping
  (`?dplyr::dplyr_by`) instead.
\end{verbatim}

\textbf{MERGING ALL DATASETS} We have: - conflict1 - disaster\_filter2 -
infant\_long - matmor\_long - neonatal\_long - under5\_long They all
have the ISO term. Hence they will be merged this way

\begin{Shaded}
\begin{Highlighting}[]
\NormalTok{full\_dataset\_list }\OtherTok{\textless{}{-}} \FunctionTok{list}\NormalTok{(matmor\_long, neonatal\_long, infant\_long, under5\_long, disaster\_filter2, conflict1)}

\NormalTok{merged\_all }\OtherTok{\textless{}{-}}\NormalTok{ full\_dataset\_list }\SpecialCharTok{\%\textgreater{}\%} \FunctionTok{reduce}\NormalTok{(full\_join, }\AttributeTok{by =} \FunctionTok{c}\NormalTok{(}\StringTok{"iso"}\NormalTok{, }\StringTok{"year"}\NormalTok{))}

\CommentTok{\#save as a CSV file and add to proccessed folder }
\FunctionTok{write.csv}\NormalTok{(merged\_all, }\StringTok{"Merged\_Data.csv"}\NormalTok{, }\AttributeTok{row.names =} \ConstantTok{FALSE}\NormalTok{)}
\end{Highlighting}
\end{Shaded}





\end{document}
